\documentclass[spanish, 12pt, a4paper]{article}

\usepackage{name}

\begin{document}

\maketitle
\thispagestyle{fancy}
\setcounter{page}{1}
\maketitle
\begin{multicols}{2}

\section{Distribución binomial}
Consideramos una población de individuos que no se reproducen y evolucionan en tiempo discreto. 
En cada paso de tiempo, cada individuo puede morir con una probabilidad $d$.

Realizamos simulaciones con diferentes probabilidades de muerte, $d = { 0.1,~0.5,~0.9 }$, 
utilizando una población inicial de $N_0 = 100$ individuos y un tiempo de simulación $t_f = 10$. 

\begin{figure}[H]
    \centering
    \includegraphics[width=0.5\textwidth]{../Ejercicio1/EvolucionProb_0.1.pdf}
    \caption{Evolución tempoal del histograma de la población con probabilidad de muerte $d = 0.1$.}
    \label{fig:binomial1}
\end{figure}

Se llevaron a cabo 1000 simulaciones para cada valor de $d$, lo que nos permite realizar análisis 
estadísticos en cada paso de tiempo. 
Las Figuras \ref{fig:binomial1}, \ref{fig:binomial2} y \ref{fig:binomial3} 
muestran los histogramas de la población en diferentes momentos del tiempo, 
comparados con la distribución binomial.

Se puede observar que rápidamente los histogramas adquieren la forma de una distribución binomial exacta y que la media se 
estabiliza en el valor esperado $d N_0$.

\begin{figure}[H]
    \centering
    \includegraphics[width=0.4\textwidth]{../Ejercicio1/EvolucionProb_0.5.pdf}
    \caption{Evolución temporal del histograma de la población con probabilidad de muerte $d = 0.5$.}
    \label{fig:binomial2}
\end{figure}

\begin{figure}[H]
    \centering
    \includegraphics[width=0.4\textwidth]{../Ejercicio1/EvolucionProb_0.9.pdf}
    \caption{Evolución temporal del histograma de la población con probabilidad de muerte $d = 0.9$.}
    \label{fig:binomial3}
\end{figure}

\section{Ecuaciones de Langevin}
Consideramos una población continua \( x(t) \), con una dinámica multiplicativa en tiempo discreto y además le añadimos un ruido aditivo. La ecuación de Langevin que describe la evolución de la población es:

\begin{equation}
    x(t+1) = a x(t) + z(t),
\end{equation}

donde la variable estocástica \( z(t) \) tiene una distribución gaussiana con media cero y desviación estándar \( \sigma \).

Simulamos doce trayectorias del sistema con la condición inicial $x(0) = 1$, utilizando $a = 1.05$ y $\sigma = 0.2$, hasta un tiempo final $t_f = 50$. 

En la Figura \ref{fig:langevinaditivo} se muestran las trayectorias simuladas. 
Se observa que las trayectorias con ruido aditivo presentan fluctuaciones alrededor de la dinámica esperada sin ruido. 
Algunas de las trayectorias logran alejarse de la esperada sin ruido debido a que el ruido afecta de manera significativa a la población 
cuando esta es pequeña, provocando su desviación de la dinámica determinista.

\begin{figure}[H]
    \centering
    \includegraphics[width=1\linewidth]{../Ejercicio2/LangevinAditivo.pdf}
    \caption{Trayectorias simuladas de la ecuación de Langevin con ruido aditivo. Las trayectorias con círculos rellenos en negro representan la dinámica exponencial sin ruido.}
    \label{fig:langevinaditivo}
\end{figure}

Se graficó (ver Figura \ref{fig:HistogramaAditivo}) la evolución de la distribución de probabilidad de las trayectorias $P(x,t)$. 
Se observa que a medida que aumenta el tiempo, la distribución presenta una mayor varianza. 
También se puede ver que, una parte de las simulaciones la población se extingue, lo que se refleja en un pequeño pico en 0.

\begin{figure}[H]
    \centering
    \includegraphics[width=1\linewidth]{../Ejercicio2/HistogramaAditivo.pdf}
    \caption{Evolución de la distribución de probabilidad $P(x,t)$ para la dinámica multiplicativa en tiempo discreto con un ruido gaussiano de media 0 y $\sigma = 0.2$.}
    \label{fig:HistogramaAditivo}
\end{figure}

Este análisis se puede realizar para la misma población, esta vez agregando un ruido multiplicativo. 
De esta forma, la dinámica se describe mediante la ecuación

\begin{equation}
    x(t + 1) = a x(t) + z(t) x(t),
\end{equation}

donde lo único que cambiamos es que el ruido ahora está multiplicado por la población en el tiempo $t$. 
En la Figura \ref{fig:langevinmultiplicativo} se pueden observar las doce de trayectorias simuladas para este caso. 
Se aprecia que las trayectorias con ruido multiplicativo tienden a modificar bruscamente la dinámica debido a que el ruido 
se amplifica con el tamaño de la población.

\begin{figure}[H]
    \centering
    \includegraphics[width=1\linewidth]{../Ejercicio2/LangevinMultiplicativo.pdf}
    \caption{Trayectorias simuladas de la ecuación de Langevin con ruido multiplicativo.
    Los circulos rellenos negros muestran la dinámica exponencial sin ruido.}
    \label{fig:langevinmultiplicativo}
\end{figure}


Se graficó además la distribución de probabilidad $P_m(x,t)$ de esta población. 
Esta se muestra en la Figura \ref{fig:HistogramaMultiplicativo}. 
Se observa que la variación es aún mayor y ocurre rápidamente al avanzar en el tiempo $t$. 
También se encuentra el pico en 0 correspondiente a las poblaciones que se extinguen debido a las fluctuaciones


\begin{figure}[H]
    \centering
    \includegraphics[width=1\linewidth]{../Ejercicio2/HistogramaMultiplicativo.pdf}
    \caption{Evolución de la distribución de probabilidad $P_m(x,t)$ para la dinámica multiplicativa con un ruido multiplicativo
    gaussiano de meida 0 y $\sigma = 0.2$.}
    \label{fig:HistogramaMultiplicativo}
\end{figure}
Como conclusión, ambos pueden ser casos posibles dentro de una población de individuos. 
Por lo tanto, la decisión de qué modelo utilizar queda sujeta a las observaciones de datos de la realidad.

\section{Simulación estocástica}
Consideramos un sistema con reproducción y competencia intraespecífica del tipo:

\begin{align*} 
    A & \xrightarrow{b} A+A \\ 
    A+A & \xrightarrow{d} A,
\end{align*}
donde $A$ representa una población con nacimientos a una tasa $b$ y muertes por competencia a una tasa $d$.
Simulamos el sistema utilizando el algorítmo de Guillespie analizando los distintos casos de las tasas
de nacimiento y muerte.

En primer lugar exploramos las dinámicas del sistema comparadas con las soluciones de la ecuación macroscópica
dada por la ecuación logistica

\begin{equation}
    \frac{dA}{dt} = bA\left(1-\frac{A}{b/d}\right)
\end{equation}

donde $b/d$ es la capacidad de carga del sistema y $b$ es la tasa de crecimiento de la población.

Veamos distintos casos según la población incial $A_0$:
\begin{enumerate}
    \item Para $A_0$ pequeño por debajo de $\frac{b/d}{2}$ la población crece como una curva cóncava
    \item Para $A_0$ grande por encima de $\frac{b/d}{2}$ y menor a $b/d$ la población crece como una curva convexa
    \item Para $A_0$ mayor a $b/d$ la población decrece hasta llegar a la capacidad de carga del sistema
\end{enumerate}

En primer lugar tomamos las tasas $b = 1$, $d = 0.01$, entonces la capacidad de carga del sistema macroscópico es $b/d = 100$.
En la figura \ref{fig:EstacionarioComparacion} se puede observar que en el estado estacionario 
la solución dada por el algoritmo de Guillespie tiene fluctuaciones alrededor de la solución dada por la ecuación logística.

\begin{figure}[H]
    \centering
    \includegraphics[width=1\linewidth]{../Ejercicio3/EstacionarioComparacion.pdf}
    \caption{Comparación de la solución estacionaria del algoritmo de Guillespie con la solución de la ecuación logística.}
    \label{fig:EstacionarioComparacion}
\end{figure}

Estudiamos además los estados transitorios para distintos valores de $b$ y $d$, las simulaciones realizadas se muestran en 
la Figura \ref{fig:Comparacion2}. Se observa un comportamiento similar al caso estacionario, 
donde la solución del algoritmo de Guillespie fluctúa alrededor de la solución de la ecuación logística.
El cambio en las tasas de nacimiento y muerte afecta la velocidad de convergencia al estado estacionario, ya que cambia
el valor de la capacidad de carga del sistema.


\begin{figure}[H]
    \centering
    \includegraphics[width=1\linewidth]{../Ejercicio3/Comparacion2.pdf}
    \caption{Comparación
    de la solución en el transitorio del algoritmo de Guillespie con la solución de la ecuación logística.}
    \label{fig:Comparacion2}
\end{figure}

Se analizó el histograma de la población en el estado estacionario para distintos valores de $b$ y $d$. 
En primer lugar se tomaron los valores $b = 0.1$ y $d = 0.001$, donde la capacidad de carga del sistema es $b/d = 100$.
Simulando 1000 evoluciones del sistema, se obtuvo el histograma de la población en el estado estacionario, el cual se muestra
en la figura \ref{fig:EstacionarioPositivoHistograma},
se observa que la población se encuentra en torno a la capacidad de carga del sistema con una distribución gaussiana con media $99.4$ (coincidente con la capacidad de carga del sistema) y
desviación estándar $9.97$.

\begin{figure}[H]
    \centering
    \includegraphics[width=1\linewidth]{../Ejercicio3/EstacionarioPositivoHistograma.pdf}
    \caption{Histograma de la población en el estado estacionario para $b = 0.1$ y $d = 0.001$.}
    \label{fig:EstacionarioPositivoHistograma}
\end{figure}

Además se analizaron los tiempos de extinción de la población para un valor de $b = 0.1$ y $d = 0.02$. Se realizaron
1000 simulaciones y se obtuvo el histograma de los tiempos de extinción de la población, el cual se muestra en la figura
\ref{fig:EstacionarioNegativoHistograma}. Se observa que la población se extingue en un tiempo finito con una media 
aproximada de $8.64$ pasos de tiempo y una desviación estándar de $6.06$.

\begin{figure}[H]
    \centering
    \includegraphics[width=1\linewidth]{../Ejercicio3/EstacionarioNegativoHistograma.pdf}
    \caption{Histograma de los tiempos de extinción de la población para $b = 0.1$ y $d = 0.02$.}
    \label{fig:EstacionarioNegativoHistograma}
\end{figure}


\end{multicols}
\end{document}